% Part: normal-modal-logic
% Chapter: axioms-systems
% Section: introduction

\documentclass[../../../include/open-logic-section]{subfiles}

\begin{document}

\olfileid[zh]{nml}{axs}{int}

\olsection{引言}

我们基于模态模型为基本模态语言建立了一种语义，并有一个关于!!a{formula}有效的概念（在所有模型的所有世界中为真），或相对于某类模型或框架有效（在该类所有模型的所有世界中为真，或基于该框架为真）的概念。逻辑学通常把这种对有效性的语义刻画与!!{derivability}的证明论概念联系起来。目标是在某个系统中定义一个!!{derivability}的概念，使得!!a{formula}是!!{derivable}当且仅当它有效。

最简单且历史上最古老的!!{derivation}系统是所谓的\zhFullName{David Hilbert}{大卫·希尔伯特}{希尔伯特}型或公理化的!!{derivation}系统。许多模态逻辑的\zhFullName{David Hilbert}{大卫·希尔伯特}{希尔伯特}型!!{derivation}系统相对容易构造：作为元理论研究（例如证明可靠性和完全性）的对象，它们很简单。然而，与自然演绎系统相比，用它们来证明!!{formula}要难得多。

在\zhFullName{David Hilbert}{大卫·希尔伯特}{希尔伯特}型!!{derivation}系统中，!!a{formula}的!!a{derivation}是一列!!{formula}，通过少数几条推理规则，从某些公理出发导出所讨论的!!{formula}。由于我们要让!!{derivation}系统与语义相匹配，就必须保证，在所有模型中，所有!!{derivable}公式的集合都为真（或在所有公理都为真的所有模型中都为真）。我们首先分离出模态逻辑为实现这一点所必需的一些性质：「正规」的模态逻辑。对于正规模态逻辑，需要假定的推理规则只有两条：分离规则和必然化。作为公理，我们取所有（代入特例）重言式，并且，视我们所处理的模态逻辑而定，取若干模态公理。即使我们只关心所有模型这一类，也必须把所有 $\Ax{K}$\iftag{notprvDiamond}{}{ and~$\Ax{Dual}$} 的代入特例都算作公理。仅凭这一点就生成了最小的正规模态逻辑~$\Log K$。

\begin{defn}
\emph{分离规则}的推理是如下推理模式
\begin{prooftree}
\AxiomC{$!A$}
\AxiomC{$!A \lif !B$}
\RightLabel{\MP}
\BinaryInfC{$!B$}
\end{prooftree}
当且仅当 $!C \ident !A \lif !B$ 时，称!!a{formula}~$!B$ 由分离规则从!!{formula} $!A$、$!C$ \emph{得出}。
\end{defn}

\begin{defn}
\emph{必然化}规则是如下推理模式
\begin{prooftree}
\AxiomC{$!A$}
\RightLabel{\Nec}
\UnaryInfC{$\Box !A$}
\end{prooftree}
当且仅当 $!B \ident \Box !A$ 时，称!!{formula}~$!B$ 由必然化从!!{formula} $!A$ 得出。
\end{defn}

\begin{defn}
一个从公理集~$\Sigma$ 出发的\emph{!!{derivation}}是一列!!{formula} $!B_1$、$!B_2$、\dots、$!B_n$，其中每个 $!B_i$ 或者是
\begin{enumerate}
\item 某个重言式的代入特例，或
\item $\Sigma$ 中某个!!a{formula}的代入特例，或
\item 由分离规则从两个!!{formula} $!B_j$、$!B_k$ 得出，其中 $j$、$k < i$，或
\item 由必然化从!!a{formula}~$!B_j$ 得出，其中 $j < i$。
\end{enumerate}
若存在这样的!!a{derivation}且 $!B_n \ident !A$，称 $!A$ \emph{从 $\Sigma$ !!{derivable}}，记为 $\Sigma \Proves !A$。
\end{defn}

有了这个定义，事实将表明：所有!!{derivable}公式的集合构成一个正规模态逻辑，并且任何!!{derivable}!!{formula}在每一个所有公理都为真的模型中都为真。!!{derivation}的这一性质称为\emph{可靠性}。其逆命题\emph{完全性}则更难证明。

\end{document}
